\section{Economic use and the remaining research task}
\label{sec:discussion121}

The role of a Neural Bellman Operator is to acquire and reuse continuation
information that changes feasible economic decisions at an acceptable total
cost. The marginal construction gives that objective a precise training and
verification account. It also exposes two costs that cannot be removed by
renaming a fitting statistic: integration must track the original risk
operator, and an exported value must retain its state-dependent level.

The acquisition budgets in this experiment change information density and
representation size within the same two-state economy. They do not establish
scalability in horizon, state dimension or innovation dimension. The retained
theorems describe those separate objects and their assumptions; their
historical experiments and unresolved regions remain in the development
companion. The stopping, accuracy-allocation and risk-information questions
therefore remain part of the NBO programme under the original economic
criteria.

A neural-specific advantage requires a distinction that a conventional
method with the same information, features, feasible policies and resources
cannot reproduce. The exactly matched hinge control makes one such boundary
explicit. Conditional adoption additionally requires the original objective
gain to pay for all acquisition and complete-service work at a stated price
and volume. These are empirical and mathematical requirements for an
economically useful NBO, not consequences of a successful optimization status
or a smaller internal loss.

The measured ReLU--hinge comparison needs a further computational distinction.
The two encodings carry identical mathematical features and readouts. Under
the frozen acquisition procedure, the coordinate-loop hinge implementation
pays for the corresponding ReLU fit and its own encoding, identity check and
audit. Those are real costs of the executed implementation, and remain in the
ledger. A conventional hinge representation can also use the same batched
evaluation and elementary operations as the ReLU implementation. Consequently
a cost difference caused by the declared copy or evaluation procedure cannot
establish a neural-specific computational advantage. Any sufficient price
region against the recorded implementations is conditional on that measured
menu. Its interpretation must also retain the other strong conventional
alternatives and the distinction between implementation cost and a limitation
of conventional representation.
